% valorile utile ale examenului din iunie 2026 (cuantile t cu grade de libertate distincte)
\providecommand{\valoriutile}{%
Valori utile: $z_{0,05}=-1{,}645$; $z_{0,025}=-1{,}960$; $z_{0,01}=-2{,}326$;
$\varphi(2{,}326)=0{,}0267$; $\ln 2 = 0{,}693$; $\sqrt{5}\approx 2{,}236$;
$\sqrt{10}\approx 3{,}162$; $\sqrt{252}\approx 15{,}87$; $\chi^2_{0,05}(1)=3{,}84$;
$\chi^2_{0,05}(2)=5{,}99$; $t^{-1}_{5}(0{,}01)=-3{,}365$; $t^{-1}_{6}(0{,}01)=-3{,}143$.}
\noindent Nume și prenume: \dotfill\ \ Grupa: \rule{2.6cm}{0.4pt}\par
\vspace{0.4ex}\noindent Data: \rule{2cm}{0.4pt}\ 2026 \hfill Timp de lucru: 2 ore\par
\vspace{0.6ex}\noindent Toate subiectele sînt obligatorii. Se acordă 1 punct din oficiu. Calculatorul este permis.\par
\vspace{0.3ex}{\small\itshape\noindent Pentru întrebările de interpretare, răspunsurile trebuie să fie concise:
maximum 3--5 fraze. La calcule se punctează atît rezultatul numeric, cît și interpretarea economică/statistică.
Unde se folosesc formule aproximative, menționați ipotezele principale.\par}
\vspace{0.5ex}{\small\noindent\valoriutile\par}
